Wir nehmen an, dass der Barbier das Messer gleichförmig und parallel zum Riemen zieht. Dann ist seine Zugkraft gleich der Gleitreibungskraft, $F = \mu_\mathrm{G}\,F_\mathrm{N}$. Für Leder auf Metall ist $\mu_\mathrm{G} = \muGO$ (FoTaBe S.~189). Das Leder drückt also mit der Normalkraft
\begin{align}
 F_\mathrm{N} &= \FNoF = \frac{\FZu}{\muG} = \FNo \approx \result{\FNoP{0}{2}}
\end{align}
gegen die Klinge.
